\chapter{Descripción general}

\begin{procedimiento}{Objetivo del capítulo}{Todos}
Presentar qué permite hacer la plataforma y cómo se relacionan sus etapas principales, sin entrar todavía en procedimientos detallados.
\end{procedimiento}

\section{Propósito de la plataforma}

La Plataforma Web Robot Explota Globos reúne en un mismo sitio las actividades principales del torneo: publicación de información para el público, registro de equipos, confirmación de asistencia, operación de competencias, comunicaciones internas y administración de datos autorizada.

Para un visitante, la plataforma funciona como punto de entrada al evento. Para un participante, permite enviar la inscripción de un robot y revisar mensajes asociados al registro. Para la organización, el sistema concentra las herramientas de preparación y operación que se explican en los capítulos posteriores.

\figuraManual{capturas/finales/MU-001-inicio-publico.png}{Página de inicio desde la cual el visitante puede consultar reglas, revisar el estado del evento o comenzar el proceso de registro.}{fig:bloque1-mu-001}

\section{Áreas principales}

\begin{tabularx}{\linewidth}{>{\bfseries}p{3.2cm}X}
\toprule
Área & Uso para el usuario\\
\midrule
Sitio público & Presenta la información general del torneo, reglas destacadas, patrocinadores y accesos al registro.\\
Registro & Permite elegir categoría, diligenciar el formulario escolar o universitario y recibir la confirmación visual del envío.\\
Asistencia & Permite confirmar participación mediante un enlace enviado por la organización.\\
Panel interno & Permite a usuarios autorizados revisar divisiones, preparar competencia y operar fases.\\
Comunicaciones & Permite gestionar plantillas, destinatarios, invitaciones, solicitudes de asistencia e historial.\\
Administración & Permite a superusuarios revisar o corregir información desde el administrador cuando la operación lo requiere.\\
\bottomrule
\end{tabularx}

\section{Relación general entre etapas}

El uso habitual comienza en el sitio público. El participante consulta la información disponible, revisa las reglas y accede al registro. Después de enviar el formulario, la organización puede revisar la información, enviar comunicaciones y solicitar confirmación de asistencia. Con los equipos confirmados, los operadores preparan la competencia, registran resultados y cierran el torneo.

\begin{nota}
Este capítulo ofrece una visión general. Los procedimientos concretos de registro se explican en el \cref{cap:registro-participantes}. Las operaciones de torneo, comunicaciones y administración se desarrollarán en los capítulos posteriores.
\end{nota}

\section{Qué no cubre este capítulo}

Este capítulo no explica cómo instalar la plataforma, configurar servidores, modificar código ni administrar bases de datos. Tampoco reemplaza las decisiones de la organización sobre calendario, reglas finales, URL de producción o responsables operativos.

\begin{resultadoesperado}
Al finalizar este capítulo, el lector debe comprender para qué sirve la plataforma, qué perfiles la usan y qué tipo de operaciones se realizan en cada área.
\end{resultadoesperado}
